% =============================================================
% packages.tex — Déclaration des paquets
% Compilation requise : XeLaTeX (fontspec charge la police système
% "Times New Roman" — installée par défaut sous Windows/MiKTeX).
% =============================================================

% --- Police et langue ---
\usepackage{fontspec}
\usepackage[french]{babel}
\usepackage{csquotes}

% --- Mise en page ---
% Recto verso (norme ESPRIT §A.1) : marges symétriques 2,5 cm imposées
% sur les quatre côtés (pas de fer à reliure asymétrique).
\usepackage[a4paper,margin=2.5cm]{geometry}
\usepackage{setspace}
\usepackage{titlesec}
\usepackage{fancyhdr}
\usepackage{lastpage}
\usepackage{chngcntr}
\usepackage{etoolbox}
% NB : ni "needspace" ni "emptypage" ne sont installés localement et la
% machine n'a pas d'accès réseau pour les récupérer (cf. formatting.tex :
% leurs effets sont reproduits avec des primitives du noyau LaTeX).

% --- Figures, tableaux, graphiques ---
\usepackage{graphicx}
\usepackage{tikz}
\usetikzlibrary{arrows.meta,positioning,shapes.geometric,calc,fit,backgrounds}
\usepackage{booktabs}
\usepackage{longtable}
\usepackage{array}
\usepackage{multirow}
\usepackage{caption}
\usepackage{subcaption}
% Type de flottant dédié « Graphique » (histogrammes, courbes), numéroté
% et listé séparément de Fig./Tab., conformément à la norme ESPRIT §C.
% Ni "float" ni "newfloat" (utilisé en interne par \DeclareCaptionType
% des versions récentes de "caption") ne sont installés hors ligne : le
% type de flottant est déclaré à la main avec les primitives du noyau
% LaTeX, exactement comme \figure et \table le sont dans classes.dtx.
% Sans compteur "within", la numérotation est continue sur tout le
% document (pas de remise à zéro par chapitre).
\makeatletter
\newcounter{graphique}
\renewcommand{\thegraphique}{\arabic{graphique}}
\def\fps@graphique{htbp}
\def\ftype@graphique{4}
\def\ext@graphique{lgr}
\def\fnum@graphique{Graph.~\thegraphique}
\newenvironment{graphique}{\@float{graphique}}{\end@float}
\newenvironment{graphique*}{\@dblfloat{graphique}}{\end@dblfloat}
\newcommand*\l@graphique{\@dottedtocline{1}{1.5em}{2.8em}}
% \listof{type}{heading} n'existe pas dans le noyau LaTeX2e : c'est en
% réalité une commande fournie par le paquet "float" (indisponible hors
% ligne). \listofgraphiques est donc récrite à la main, sur le modèle
% exact de \listoffigures dans classes.dtx.
\newcommand{\listgraphiquename}{Liste des graphiques}
\newcommand\listofgraphiques{%
  \if@twocolumn
    \@restonecoltrue\onecolumn
  \else
    \@restonecolfalse
  \fi
  \chapter*{\listgraphiquename}%
  \@mkboth{\MakeUppercase\listgraphiquename}{\MakeUppercase\listgraphiquename}%
  \addcontentsline{toc}{chapter}{\listgraphiquename}%
  \@starttoc{lgr}%
  \if@restonecol\twocolumn\fi
}
% Les classes standard n'ajoutent pas les listes des figures et des
% tableaux au sommaire. Les deux redéfinitions ci-dessous les y ajoutent
% au bon numéro de page, comme la liste des graphiques ci-dessus.
\renewcommand\listoffigures{%
  \if@twocolumn
    \@restonecoltrue\onecolumn
  \else
    \@restonecolfalse
  \fi
  \chapter*{\listfigurename}%
  \@mkboth{\MakeUppercase\listfigurename}{\MakeUppercase\listfigurename}%
  \addcontentsline{toc}{chapter}{\listfigurename}%
  \@starttoc{lof}%
  \if@restonecol\twocolumn\fi
}
\renewcommand\listoftables{%
  \if@twocolumn
    \@restonecoltrue\onecolumn
  \else
    \@restonecolfalse
  \fi
  \chapter*{\listtablename}%
  \@mkboth{\MakeUppercase\listtablename}{\MakeUppercase\listtablename}%
  \addcontentsline{toc}{chapter}{\listtablename}%
  \@starttoc{lot}%
  \if@restonecol\twocolumn\fi
}
\makeatother

% --- Encadrés (principes clés, alertes, éléments en attente) ---
\usepackage{xcolor}
\usepackage{tcolorbox}
\tcbuselibrary{skins,breakable}

% --- Bibliographie ---
\usepackage[
  backend=biber,
  style=numeric,
  sorting=nyt,
  maxbibnames=99,
  urldate=long
]{biblatex}

% --- Liens et références croisées (chargés en dernier) ---
% Bibliographie en drapeau : evite l'etirement des URL par la
% justification et quatre depassements de ligne.
\appto\bibsetup{\raggedright}

% NB : "microtype" n'est pas installé localement et la machine n'a pas d'accès
% réseau pour le récupérer (même situation que "needspace" et "emptypage",
% cf. plus haut ; vérifié le 30/08/2026). Les débordements de ligne sont donc
% traités par la réserve d'étirement \emergencystretch fixée dans
% formatting.tex, qui les ramène tous sous le seuil de 10 pt.
\usepackage{hyperref}
\usepackage{cleveref}
